\section{Analytic Tasks and the Changing Locus of Agency}
\label{sec:tasks}

We compare systems by the analytical work scientists seek to accomplish and the computational contributions used to support it. For each task, we examine what computation performs, which decisions scientists supply or retain, and how they can inspect and revise the result. The Assistance $\times$ Visualization Modality framework organizes the mechanisms and environments involved; the locus of analytic agency identifies differences in initiative and control over particular actions. These comparisons connect system capabilities to the evaluation questions needed to assess their analytical benefits.

The five task categories introduced in Section~\ref{sec:foundations} describe related purposes rather than successive stages. Scientists revisit selections, comparisons, and explanations as questions develop. We therefore compare systems around consequential actions, reserving implementation detail for Section~\ref{sec:assistance} and environmental conditions for Section~\ref{sec:modality}. An interaction can support several tasks, but its presence does not establish that every task is supported equally well.

The task symbols in Figure~\ref{fig:assistance-modality-map} allow these purposes to be followed across assistance modes and environments; sharing a symbol does not imply equivalent task support or performance.

\subsection{Navigation and Multiscale Orientation}
\label{sec:tasks-navigation}

Navigation concerns locating relevant information while retaining enough context to understand its relationships. In life-science analysis, the problem includes movement between a local structure and its wider network, or between an observed specimen and a computed feature space. Panning and zooming change the visible region; computational organization also changes the structure through which the scientist navigates. The distinction matters because an accessible overview is already a selection of what to emphasize.

Compressed Adjacency Matrices organize regulatory networks so that motifs, paths, and neighborhoods can be inspected through highlighting, filtering, and selected views.~\cite{AnchorCompressedMatrices} Facetto connects tissue images with feature distributions, projections, and phenotype organization.~\cite{AnchorFacetto} The former makes network structure traversable through a computed representation; the latter supports inspection across spatial and derived descriptions of cells. These examples concern orientation among relationships and representations. They should not be recast as demonstrated autonomous navigation or guided scale changes without evidence of those mechanisms.

Agency lies partly in who chooses the next object of attention and partly in how the available overview was constructed. A scientist may control every movement while relying on computational abstractions to determine what is readily visible. Evaluation should consequently distinguish reaching a target from maintaining context, recognizing omitted structure, and returning to an earlier observation. A display may make traversal easier without improving those other outcomes. The modality comparison must identify which of these costs a change of environment actually addresses.

\subsection{Selection, Filtering, and Precision Interaction}
\label{sec:tasks-selection}

Selection establishes which entities or relationships an analysis concerns. Direct pointing specifies an object; filtering specifies conditions; assisted selection may infer membership from examples or translate a request into a set. These actions can produce similar-looking highlights while assigning different responsibility for the selected population. Precision therefore has both an interaction meaning, selecting the intended object, and an analytical meaning, forming a set appropriate to the question.

Facetto makes this distinction consequential when experts select image regions or feature groups and revise phenotype labels. Feedback can affect later classifier assignments beyond the cells directly inspected.~\cite{AnchorFacetto} In FathomGPT, a language request can instead determine the queried records through generated SQL.~\cite{AnchorFathomGPT} The scientist provides examples in one case and a description in the other; computation contributes to the resulting membership. Checking a selected set requires access to the relevant basis for membership, not merely confirmation that a selection action succeeded.

NUI-VR$^2$ illustrates a spatial counterpart: user-supplied seeds and parameters influence computed segmentation boundaries.~\cite{AnchorNUIVR} A precise gesture and an accurate segmentation are related but distinct accomplishments. Its technical segmentation measures assess the resulting boundary, not the full effort of choosing seeds, detecting errors, and correcting them. Across these mechanisms, the useful comparison is how scientists specify intent, inspect the selected result, and recover from unintended inclusion or omission. That comparison can then inform whether planar input, language, or spatial interaction suits the particular selection problem.

\subsection{Comparison and Differentiation}
\label{sec:tasks-comparison}

Comparison requires a reference: another observation, condition, model, source, or expected pattern. Computation can establish correspondence, derive differences, and organize cases for inspection, but those choices shape what a visible difference means. A salient contrast may reflect measurement, representation, or modeling choices, rather than a biological distinction. The analyst's responsibility therefore includes deciding which comparisons are meaningful and which alternative explanations remain plausible.

RetainVis and MediSyn locate this work at different evidential levels. RetainVis exposes risk predictions and feature contributions across patients and cohorts, with editing and retraining that can change inputs or model state.~\cite{AnchorRetainVis} MediSyn aligns biomedical assertions across datasets and makes provenance, coverage, and disagreement inspectable.~\cite{AnchorMediSyn} In the former, comparison can investigate how a computed inference responds to intervention; in the latter, it can investigate where sources support or contradict an interpretation. An input edit is not evidence of a real-world causal effect, just as agreement among displayed assertions is not by itself proof of their independence or correctness.

The distinction also changes the evaluation target: RetainVis's model behavior and MediSyn's evidence-selection performance are not interchangeable outcomes. Section~\ref{sec:evaluation-effectiveness} separates computational behavior from analytical and scientific benefit.

Comparison can also interrogate the representation itself. focusedMDS/distnet exposes discrepancies between projected and high-dimensional distances, while PathwayMatrix uses ordering, grouping, and lenses to inspect pathway relationships.~\cite{ExpandedE34,ExpandedE37} PAE Viewer links predicted alignment error with protein structure and cross-links, making model reliability an object of inspection rather than an implicit property of the display.~\cite{ExpandedE49} By contrast, CoCoNUT's sequence comparisons and Toxygates' clustering and enrichment organize evidence about genomic similarity and response patterns.~\cite{ExpandedE35,ExpandedE36} These mechanisms support different comparisons even when their interfaces share planar views.

\subsection{Sensemaking and Hypothesis Development}
\label{sec:tasks-sensemaking}

Sensemaking connects observations with possible explanations and questions for further investigation. Assistance may expose a pattern, suggest a candidate, or execute an analysis that tests part of an interpretation. The scientific work includes deciding what the result warrants and what further evidence is needed. A computationally generated possibility becomes useful through that relationship to an investigation, rather than through generation alone.

RNA-SeqEZPZ provides interactive outputs from a configured transcriptomics pipeline, while FathomGPT lets scientists formulate and refine requests that become queries and visual results.~\cite{AnchorRNASeqEZPZ,AnchorFathomGPT} Both support inquiry, but the route from question to computation differs. In the first, scientists configure an established sequence and interpret its output; in the second, they also inspect the translation of language into analytic operations. Neither route transfers responsibility for biological interpretation merely because execution is delegated.

Metis and ProteinSketch extend the task toward constructing candidates: feedback and spatial specifications shape which molecular possibilities are generated.~\cite{AnchorMetis,AnchorProteinSketch} Their mechanisms are compared in Sections~\ref{sec:assistance-adaptive} and~\ref{sec:assistance-immersive}. Here, the task consequence is that scientists must assess not only a candidate, but also the constraints and omissions that shaped the candidate space.

The agency consequence is a change in the space of possibilities presented to the scientist. Exposing why a candidate was generated, what constraints shaped it, and what alternatives were not considered can support critical interpretation. These are design implications of the workflows, not evidence that either system has demonstrated unbiased exploration. Section~\ref{sec:evaluation} considers how to assess those claims.

\subsection{Coordination and Collaborative Reasoning}
\label{sec:tasks-coordination}

Collaborative analysis requires more than making the same visualization available to several people. Participants must divide work, communicate interpretations, recognize changes, and determine which results or revisions can be accepted. Computational assistance can organize evidence for discussion or coordinate analytical actions. The distinction identifies whether a system primarily supports shared interpretation or the management of interdependent work.

PeCaX brings computed variant annotations and pathway relationships into a workflow intended for molecular-tumor-board inspection.~\cite{AnchorPeCaX} Its relevance to collective interpretation does not, by itself, establish a synchronous multiuser interface or automated management of disagreement. NeuroBlocks more explicitly connects segmentation and proofreading with task assignment, approval, provenance, and restoration.~\cite{AnchorNeuroBlocks} Here, coordination includes who acts on a result, how its status changes, and how prior work can be audited or recovered.

These functions distribute authority as well as effort. A proposed correction, an accepted segmentation, and a communicated interpretation need not carry the same status or belong to the same participant. Group-visible records of such distinctions can support common ground; a shared display alone cannot supply them. PeCaX's processing results and NeuroBlocks's qualitative collaborations establish different aspects of the implemented workflows. Claims about improved collective reasoning additionally require evidence about participants' understanding, reconciliation of interpretations, and use of the resulting decisions.

\subsection{Relationships Across Tasks}
\label{sec:tasks-synthesis}

Tasks connect across a workflow, so an intervention's consequences may emerge after the action it directly supports. Consider an illustrative analysis: a scientist locates a tissue region, selects a cell population, compares its characteristics, proposes an explanation, and discusses it with collaborators. If assistance changes the selection, it also changes the basis of the comparison and the evidence supporting the explanation. This sequence is an analytical illustration, not a claim that one selected system implements every step.

Table~\ref{tab:task-agency-synthesis} summarizes the resulting comparison questions. For each task, the comparison identifies the computational contribution, the decisions scientists make or delegate, and the opportunities to inspect and revise results. These differences connect the mechanisms examined next to specific evaluation requirements: preserving context, forming appropriate selections, establishing meaningful comparisons, assessing interpretations, and coordinating shared decisions.

% Requires booktabs and tabularx. Questions synthesize the discussion, not measured outcomes.
\begin{table*}[t]
\centering
\caption{Task-centered questions for comparing assistance mechanisms. The final column identifies evaluation questions; it does not claim that every cited system has evaluated them.}
\label{tab:task-agency-synthesis}
\small
\setlength{\tabcolsep}{4pt}
\renewcommand{\arraystretch}{1.15}
\begin{tabularx}{\textwidth}{@{}>{\raggedright\arraybackslash}p{.15\textwidth}>{\raggedright\arraybackslash}X >{\raggedright\arraybackslash}X >{\raggedright\arraybackslash}X@{}}
\toprule
Task & Computational participation & Consequential human choice & Evaluation question \\
\midrule
Navigation and orientation & Organize structure; connect spatial and derived views. & Choose attention and interpret the context retained or omitted. & Can analysts locate relevant structure, retain context, and return to prior observations? \\
Selection and precision & Infer membership, translate set descriptions, or compute boundaries from seeds. & Specify intent and correct inclusion or omission. & Does the result match the analytical question, and how difficult is correction? \\
Comparison and differentiation & Align evidence or expose differences in predictions and contributions. & Choose reference conditions and interpret the source of a difference. & Can analysts distinguish source disagreement, model behavior, and biological variation? \\
Sensemaking and hypothesis development & Execute inquiries or generate candidates under feedback and constraints. & Formulate questions, inspect assumptions, and determine what warrants further testing. & Are interpretations traceable, alternatives inspectable, and claims supported by appropriate evidence? \\
Coordination and collaborative reasoning & Organize shared evidence or track assignments, approvals, and revisions. & Divide work, authorize changes, and reconcile interpretations. & Do participants understand the status, provenance, and consequences of one another's actions? \\
\bottomrule
\end{tabularx}
\end{table*}

